\chapter{ANALYSE DES BESOINS}

\section{Introduction}

La phase d'analyse des besoins constitue le socle sur lequel repose
l'intégralité du processus de conception et de développement. Elle vise
à recenser, formaliser et hiérarchiser l'ensemble des exigences
auxquelles le système devra répondre, en tenant compte des contraintes
organisationnelles propres à l'environnement d'accueil.

Dans le contexte de ce projet, cette analyse prend une dimension
particulière~: la plateforme est destinée à des acteurs aux profils
et aux périmètres d'action hétérogènes, répartis selon une hiérarchie
fonctionnelle stricte. Il est donc indispensable d'identifier avec
précision les responsabilités de chaque intervenant, les
fonctionnalités qui leur sont associées, ainsi que les contraintes
non fonctionnelles qui encadrent l'ensemble du système.

Ce chapitre s'articule autour de quatre axes~: l'identification des
acteurs et de leurs rôles, la formalisation des besoins fonctionnels
et non fonctionnels, la modélisation des cas d'utilisation sous forme
de diagrammes UML, et enfin la traduction de l'ensemble en un
backlog produit Scrum réparti sur six sprints.

\section{Identification des acteurs}

\subsection{Définition}

Un acteur, au sens de la modélisation UML, désigne toute entité
extérieure au système qui interagit avec lui, qu'il s'agisse d'un
être humain, d'un autre système logiciel ou d'un processus
automatisé. Dans le cadre de ce projet, les acteurs se déclinent en
deux catégories~: les acteurs humains, représentant les différents
profils d'utilisateurs de la plateforme, et un acteur système, le
service de prédiction FastAPI, qui intervient de manière automatisée
dans les flux d'analyse.

\subsection{Acteurs humains}

La plateforme est conçue pour être utilisée par cinq profils
distincts, organisés selon une hiérarchie fonctionnelle à trois
niveaux. Cette organisation conditionne directement le périmètre
d'accès de chaque utilisateur aux données et aux fonctionnalités du
système.

\begin{table}[htbp]
    \centering
    \caption{Acteurs humains et périmètres d'action}
    \label{tab:acteurs}
    \begin{tabular}{|p{3.2cm}|p{4.5cm}|p{6cm}|}
        \hline
        \textbf{Acteur} & \textbf{Niveau hiérarchique}
        & \textbf{Périmètre d'action} \\
        \hline
        Super Admin
        & Niveau 1~: administration globale
        & Accès complet à l'ensemble du système~: gestion des
          utilisateurs, réinitialisation de la base de données,
          supervision multi-agences, vidage de tables. \\
        \hline
        Admin Ville
        & Niveau 2~: administration territoriale
        & Validation des inscriptions, supervision des agences
          rattachées à sa ville, consultation des tableaux de bord
          de son périmètre géographique. \\
        \hline
        Chef d'Agence
        & Niveau 2~: supervision opérationnelle
        & Validation des inscriptions des comptes de type 'agent marketing' et 'agent commercial'
          rattachées à son agence, lancement
d’analyses, génération de don-
nées mock, validation ou rejet des recommandations soumises par les
          agents de son agence, consultation des fiches clients et
          des KPIs de son agence, accès au tableau de bord. \\
        \hline
        Agent Commercial
        & Niveau 3~: exécution terrain
        & Consultation des fiches clients, création de recommandations de type 'commercial', validation ou rejet des recommandations de type 'commercial'. \\
        \hline
        Agent Marketing
        & Niveau 3~: exécution terrain
        & Consultation des fiches clients, création de recommandations de type 'marketing', validation ou rejet des recommandations de type 'marketing'.\\
        \hline
    \end{tabular}
\end{table}

\subsection{Acteur système}

En complément des acteurs humains, le service \textbf{FastAPI}
constitue un acteur système à part entière. Opérant sur le
port~8000, ce service expose les endpoints d'inférence ML et est
appelé automatiquement par l'application Django lors du lancement
d'une analyse. Son rôle est d'exécuter les prédictions de churn sur
l'ensemble du portefeuille clients et de retourner les résultats
sous forme structurée.

Sa particularité réside dans le fait que son indisponibilité ne
bloque pas le fonctionnement de la plateforme~: un mécanisme de
\textit{fallback} local, s'appuyant sur le pipeline ML chargé
directement depuis le fichier \texttt{churn\_model\_v1.pkl}, prend
automatiquement le relais afin d'assurer la continuité de service.

\section{Besoins fonctionnels}

Les besoins fonctionnels décrivent l'ensemble des fonctionnalités
que le système doit offrir pour satisfaire les objectifs métier
identifiés. Ils sont regroupés par module fonctionnel, chacun
correspondant à une couche logicielle distincte de l'architecture.

\subsection{Module Authentification et Gestion des Comptes}

Ce module constitue le point d'entrée obligatoire de l'ensemble des
utilisateurs. Il garantit la sécurité des accès et orchestre le
cycle de vie des comptes au sein de la hiérarchie fonctionnelle.

\begin{itemize}

    \item \textbf{BF-01~— Inscription avec validation
    hiérarchique}~: tout nouvel utilisateur doit pouvoir soumettre
    une demande d'inscription en renseignant son identité, son rôle
    souhaité et son agence de rattachement. Le compte est créé avec
    un statut \textit{en attente} et ne devient actif qu'après
    validation explicite par le supérieur hiérarchique compétent~:
    un \textit{admin\_ville} pour les chefs d'agence, un
    \textit{chef\_agence} pour les agents. Cette validation
    s'effectue par email, ce qui implique la transmission
    automatique d'une notification au supérieur dès la soumission
    du formulaire.

    \item \textbf{BF-02~— Authentification à double facteur}~:
    l'accès à la plateforme repose sur un mécanisme
    d'authentification en deux étapes successives. La première
    consiste en la vérification classique des identifiants (nom
    d'utilisateur et mot de passe). En cas de succès, un code OTP à
    six chiffres est généré, doté d'une durée de validité de dix
    minutes, et transmis à l'adresse email de l'utilisateur via le
    service SendGrid. La session n'est ouverte qu'après saisie et
    validation de ce code sur l'écran dédié.

    \item \textbf{BF-02.1~— Gestion des sessions et déconnexion}~:
    l'utilisateur doit pouvoir se déconnecter explicitement de la
    plateforme via un bouton dédié, ce qui entraîne la destruction
    immédiate de sa session côté serveur. Les sessions expirées
    (inactivité prolongée ou expiration du token OTP) doivent être
    automatiquement invalidées. Un mécanisme de timeout de session
    doit être implémenté pour déconnecter automatiquement les
    utilisateurs inactifs après une période d'inactivité configurable.

    \item \textbf{BF-03~— Journalisation des connexions}~: chaque
    connexion réussie doit être enregistrée dans la base de données
    avec l'horodatage, l'adresse IP de la machine cliente et le
    \textit{user-agent} du navigateur utilisé. Ces données
    constituent un journal d'audit permettant de détecter toute
    activité anormale.

    \item \textbf{BF-04~— Réinitialisation du mot de passe}~: en
    cas d'oubli, l'utilisateur doit pouvoir déclencher une
    procédure de réinitialisation par email. Un lien sécurisé à
    usage unique, assorti d'une durée de validité limitée, lui est
    transmis pour lui permettre de définir un nouveau mot de passe.

    \item \textbf{BF-05~— Gestion des comptes par
    l'administration}~: les administrateurs disposent d'une
    interface dédiée permettant d'activer, de désactiver ou de
    supprimer un compte utilisateur, ainsi que de modifier son rôle
    ou son agence de rattachement.

    \item \textbf{BF-06~— Filtrage dynamique des rôles}~: lors de
    l'inscription, les rôles proposés à l'utilisateur peuvent être
    restreints dynamiquement selon la hiérarchie~: seuls les rôles
    inférieurs ou équivalents à celui de l'utilisateur connecté
    peuvent être attribués.

    \item \textbf{BF-07~— Affectation d'agence}~: un formulaire
    dédié permet d'associer un utilisateur à une agence valide
    référencée en base. Cette affectation conditionne directement
    la portée des données auxquelles l'utilisateur aura accès dans
    l'ensemble de la plateforme.

\end{itemize}

\subsection{Module Gestion des Données Clients}

Ce module gère l'ensemble du cycle de vie des données clients au
sein de la plateforme, depuis l'import du jeu de données jusqu'au
stockage structuré des résultats de prédiction.

\begin{itemize}

    \item \textbf{BF-08~— Import multi-formats}~: la plateforme
    permet actuellement l'import de datasets clients au format CSV.
    L'architecture du module d'import a cependant été conçue de
    manière extensible afin d'autoriser l'ajout futur de formats
    supplémentaires tels qu'Excel (\texttt{.xls}, \texttt{.xlsx}),
    TXT ou TSV.

    \item \textbf{BF-09~— Validation de la structure des
    données}~: avant tout traitement, le fichier importé doit être
    soumis à une validation automatique vérifiant la présence des
    colonnes obligatoires, l'absence de doublons sur l'identifiant
    client et la conformité des types de données. Toute anomalie
    détectée est restituée à l'utilisateur sous forme de message
    d'erreur explicite.

    \item \textbf{BF-10~— Stockage du profil client complet}~:
    chaque client est représenté par un enregistrement
    \texttt{ClientChurn} centralisant l'ensemble de ses variables
    comportementales et contractuelles pertinentes pour la
    prédiction du churn, ainsi que les résultats calculés~: score
    de churn et prédiction binaire.

    \item \textbf{BF-11~— Génération de données simulées}~: un
    générateur de données fictives (\texttt{core/mock\_data.py})
    doit permettre de créer un dataset de clients simulés, dont les
    caractéristiques respectent des distributions statistiques
    réalistes. Cette fonctionnalité est accessible via l'endpoint
    \texttt{/generer-mock/} et suit le même pipeline d'analyse que
    les données réelles.

    \item \textbf{BF-12~— Isolation des données par agence}~: les
    données clients doivent être strictement isolées par agence. Un
    utilisateur ne doit en aucun cas pouvoir consulter, modifier ou
    analyser les données appartenant à une agence autre que la
    sienne. Cette isolation est assurée à tous les niveaux~:
    requêtes Django, vues et interface utilisateur.

\end{itemize}

\subsection{Module Analyse et Prédiction}

Ce module constitue le cœur fonctionnel de la plateforme. Il
orchestre l'ensemble du processus d'analyse, depuis le déclenchement
par l'utilisateur jusqu'à la restitution des scores et des métriques.

\begin{itemize}

    \item \textbf{BF-13~— Lancement d'analyse par
    l'utilisateur}~: l'agent marketing ou le super admin doit
    pouvoir déclencher une analyse depuis l'interface web
    (\texttt{/lancer-analyse/}), en choisissant entre l'import d'un
    fichier CSV ou l'utilisation de données mock. L'analyse porte
    sur l'ensemble du portefeuille clients de l'agence.

    \item \textbf{BF-14~— Appel prioritaire au service
    FastAPI}~: lors de chaque analyse, l'application Django doit
    en priorité solliciter le service FastAPI via l'endpoint
    \texttt{/api/predict/batch}. Ce service reçoit les données au
    format CSV, exécute les prédictions sur l'ensemble du
    portefeuille et retourne les résultats sous forme JSON,
    comprenant notamment les scores individuels, les agrégats par
    niveau de risque et le taux de churn global.

    \item \textbf{BF-15~— Fallback sur le pipeline ML local}~: si
    le service FastAPI est indisponible au moment du lancement de
    l'analyse, l'application doit basculer automatiquement sur le
    pipeline ML local (\texttt{core/ml\_service.py}), qui charge le
    modèle depuis le fichier \texttt{.pkl} et produit les scores de
    manière transparente pour l'utilisateur, sans message d'erreur
    bloquant.

    \item \textbf{BF-16~— Mise à jour des scores clients}~: à
    l'issue de l'analyse, les champs \texttt{score\_churn} et
    \texttt{churn\_predit} de chaque enregistrement
    \texttt{ClientChurn} doivent être mis à jour en base de
    données.

    \item \textbf{BF-17~— Traçabilité des analyses}~: chaque
    session d'analyse doit être enregistrée dans un enregistrement
    \texttt{AnalyseSession} associant l'utilisateur déclencheur,
    l'horodatage, le type d'analyse et les métriques globales
    retournées. L'historique de ces sessions est consultable et
    archivable depuis l'interface.

    \item \textbf{BF-18~— Explications SHAP par client}~: pour
    chaque client, la plateforme doit être en mesure de fournir une
    explication de sa prédiction en termes de contribution de
    chaque variable au score calculé, via la méthode SHAP
    (\textit{SHapley Additive exPlanations}). Ces explications sont
    stockées via le modèle \texttt{ShapValeur} et peuvent être
    exploitées pour enrichir l'interprétabilité des prédictions au
    niveau de la fiche client.

\end{itemize}

\subsection{Module Tableau de Bord et Interface de Visualisation}

Ce module regroupe l'ensemble des fonctionnalités d'affichage, de
consultation et d'interaction offertes aux différents profils
d'utilisateurs.

\begin{itemize}

    \item \textbf{BF-19~— Tableau de bord avec KPIs}~: la page
    d'accueil de la plateforme présente les indicateurs clés de
    performance calculés à partir de la dernière analyse~: taux de
    churn global, nombre de clients à risque élevé, nombre de
    recommandations en attente de validation, et nombre de
    notifications non lues. \textbf{Cette fonctionnalité est
    réservée aux chefs d'agence, administrateurs de ville et super
    administrateurs. Les agents (commercial et marketing) n'ont pas
    accès au tableau de bord KPIs et sont redirigés vers une page
    d'accueil restreinte.}

    \item \textbf{BF-19.1~— Tableau de bord multi-agences pour
    admin\_ville}~: l'administrateur de ville doit pouvoir accéder à
    une vue consolidée présentant les agrégats de l'ensemble des
    agences rattachées à sa ville (taux de churn par agence, nombre
    de clients à risque par agence, comparatifs inter-agences). Les
    droits différenciés entre chef\_agence (vue limitée à son agence)
    et admin\_ville (vue multi-agences) doivent être explicitement
    formalisés. Le super admin dispose quant à lui d'une vue globale
    sur l'ensemble du système.

    \item \textbf{BF-20~— Liste paginée des clients avec
    filtres et tri}~: la plateforme doit proposer une vue tabulaire des
    clients, paginée et filtrée selon les résultats de prédiction
    et les critères de recherche disponibles. Les critères de tri
    doivent être formalisés~: tri par score de churn (croissant ou
    décroissant), tri par date de création, tri par nom du client,
    et tri par niveau de risque (élevé/faible). L'affichage est
    limité aux clients appartenant à l'agence de l'utilisateur
    connecté.

    \item \textbf{BF-21~— Fiche client détaillée}~: chaque client
    doit disposer d'une fiche individuelle regroupant ses données
    de profil, sa prédiction de churn, les informations
    d'interprétation SHAP illustrant les facteurs contributifs à sa
    prédiction, et l'historique des recommandations qui lui sont
    associées.

    \item \textbf{BF-21.1~— Consultation de l'historique des
    recommandations par client}~: depuis la fiche client, l'utilisateur
    doit pouvoir consulter l'historique complet des recommandations
    associées à ce client, avec pour chaque recommandation~: le
    type (commercial ou marketing), le statut courant, la date de
    création, l'auteur de la recommandation, et le contenu. Cette
    fonctionnalité permet de suivre l'évolution des actions de
    rétention entreprises pour un client donné.

    \item \textbf{BF-22~— Export PDF de la fiche client}~: depuis
    la fiche client, l'utilisateur doit pouvoir générer et
    télécharger un document PDF mettant en forme l'ensemble des
    informations du client, y compris le score de churn. La
    génération est assurée par l'outil \texttt{wkhtmltopdf} via la
    bibliothèque \texttt{pdfkit}.

    \item \textbf{BF-23~— Réinitialisation du tableau de bord}~:
    un mécanisme de réinitialisation (\texttt{/reinitialiser/}) doit
    permettre au chef d'agence de vider l'ensemble des données
    clients, recommandations et sessions d'analyse associées à son
    agence, afin de repartir d'un état initial pour une nouvelle
    campagne d'analyse.

\end{itemize}

\subsection{Module Recommandations}

Ce module formalise le processus de suivi des actions de rétention
client, depuis leur création par les agents jusqu'à leur validation
par la hiérarchie.

\begin{itemize}

    \item \textbf{BF-24~— Création de recommandation}~: un agent
    commercial doit pouvoir, depuis la fiche d'un client identifié
    comme étant à risque, rédiger et soumettre une recommandation
    d'action commerciale. Celle-ci est enregistrée avec le statut
    \textit{en attente} et déclenche immédiatement une notification
    à destination du chef d'agence.

    \item \textbf{BF-25~— Validation hiérarchique}~: le chef
    d'agence reçoit une notification l'invitant à consulter les
    recommandations en attente. Il dispose de deux options~:
    valider la recommandation, ce qui la fait passer au statut
    \textit{validée}, ou la rejeter en motivant sa décision, ce qui
    génère un enregistrement \texttt{RejetRecommandation} et
    notifie l'agent de l'issue.

    \item \textbf{BF-26~— Suivi du cycle de vie}~: chaque
    recommandation doit pouvoir être suivie tout au long de son
    cycle de vie à travers quatre statuts successifs~:
    \textit{pending}, \textit{validated}, \textit{rejected} et
    \textit{completed}. L'horodatage de chaque transition est
    conservé en base.

\end{itemize}

\subsection{Module Notifications}

\begin{itemize}

    \item \textbf{BF-27~— Génération automatique de
    notifications}~: le système doit générer automatiquement des
    notifications in-app à destination des utilisateurs concernés
    lors des événements métier suivants~: soumission d'une nouvelle
    recommandation, validation requise, rejet d'une recommandation
    et alerte churn pour un client à risque élevé.

    \item \textbf{BF-28~— Notification par email via SendGrid}~:
    les notifications par email sont exclusivement réservées aux
    événements liés au compte et à l'authentification~: création de
    compte (validation hiérarchique), authentification OTP,
    réinitialisation de mot de passe, et journalisation des
    connexions. Les notifications liées aux analyses, aux
    recommandations et aux alertes churn ne déclenchent PAS
    d'envoi d'email, seule la notification in-app est utilisée.
    L'envoi est effectué par \texttt{core/otp\_email\_service.py}
    avec un gabarit HTML personnalisé, depuis l'adresse définie
    dans la variable d'environnement \texttt{DEFAULT\_FROM\_EMAIL}.

    \item \textbf{BF-29~— Gestion des notifications par
    l'utilisateur}~: chaque utilisateur doit pouvoir consulter la
    liste de ses notifications, les marquer comme lues
    individuellement ou en masse, les archiver et les supprimer. Un
    compteur de notifications non lues est affiché en permanence
    dans l'interface.

\end{itemize}

\section{Besoins non fonctionnels}

Les besoins non fonctionnels définissent les contraintes de qualité,
de performance et de sécurité auxquelles le système doit se
conformer, indépendamment des fonctionnalités qu'il offre. Ils
conditionnent directement les choix d'architecture et
d'implémentation.

\subsection{Performance}

\begin{itemize}

    \item \textbf{BNF-01}~: les prédictions individuelles affichées
    sur la fiche client doivent être restituées avec un temps de
    réponse compatible avec une utilisation interactive fluide de
    la plateforme, que le calcul soit effectué par le service
    FastAPI ou par le pipeline local.

    \item \textbf{BNF-02}~: en cas d'indisponibilité du service
    FastAPI, le basculement sur le pipeline ML local doit être
    transparent pour l'utilisateur et ne doit induire aucune
    interruption de service ni message d'erreur bloquant.

    \item \textbf{BNF-03}~: le tableau de bord principal doit se
    charger en moins de trois secondes pour un portefeuille de
    taille standard, ce qui implique une optimisation des requêtes
    SQL et l'utilisation de \texttt{select\_related} sur les
    relations de clés étrangères.

\end{itemize}

\subsection{Sécurité}

\begin{itemize}

    \item \textbf{BNF-04}~: l'ensemble des pages de la plateforme
    doit être protégé par le décorateur \texttt{@login\_required}
    de Django. Toute tentative d'accès sans session active doit
    rediriger l'utilisateur vers la page de connexion.

    \item \textbf{BNF-05}~: le mécanisme OTP doit garantir
    l'unicité d'usage de chaque code~: un code utilisé ne peut plus
    être réemployé, même s'il n'a pas encore expiré. Le service
    \texttt{core/otp\_service.py} gère les statuts
    \texttt{ACTIVE}, \texttt{USED}, \texttt{EXPIRED} et
    \texttt{REVOKED}, ainsi que la limitation du nombre de
    tentatives (trois tentatives maximum).

    \item \textbf{BNF-06}~: les endpoints sensibles de
    l'application doivent être protégés contre les attaques par
    force brute grâce à un mécanisme de limitation du nombre de
    requêtes par adresse IP, prévu pour être intégré au niveau des
    endpoints sensibles.

    \item \textbf{BNF-07}~: les mots de passe doivent être stockés
    exclusivement sous forme de hachage selon l'algorithme PBKDF2
    par défaut de Django. Aucun mot de passe en clair ne doit
    transiter dans les logs applicatifs.

    \item \textbf{BNF-08}~: l'isolation des données par agence doit
    être appliquée à tous les niveaux de la pile applicative~: au
    niveau des requêtes ORM (\texttt{filter(agence=...)}), au
    niveau des vues Django, et au niveau de l'interface
    utilisateur. Une violation de ce principe constituerait une
    faille de sécurité majeure.

    \item \textbf{BNF-09}~: l'ensemble des variables sensibles
    (clé secrète Django, identifiants PostgreSQL, clé API SendGrid)
    doit être externalisé dans un fichier \texttt{.env} chargé via
    \texttt{python-decouple}. Ces variables ne doivent jamais
    apparaître en clair dans le code source.

\end{itemize}

\subsection{Maintenabilité et évolutivité}

\begin{itemize}

    \item \textbf{BNF-10}~: l'architecture modulaire du projet,
    organisée en quatre applications Django distinctes
    (\texttt{accounts}, \texttt{core}, \texttt{dashboard},
    \texttt{learning}), doit garantir l'indépendance des
    composants. Une modification dans le module \texttt{learning}
    ne doit pas nécessiter de changement dans \texttt{accounts}, et
    vice versa.

    \item \textbf{BNF-11}~: le service FastAPI doit pouvoir être
    mis à jour, remplacé ou redéployé indépendamment de
    l'application Django, sans modification du code métier. Le
    contrat d'interface entre les deux services est formalisé par
    les schémas Pydantic définis dans
    \texttt{churn\_pfe/backend/app.py}.

    \item \textbf{BNF-12}~: le code source doit être rédigé en
    français, commenté et documenté au moyen de \textit{docstrings}
    pour les fonctions critiques. Les messages de commit Git
    doivent respecter la convention \texttt{[module] Description
    courte de la modification}.

\end{itemize}

\subsection{Compatibilité et portabilité}

\begin{itemize}

    \item \textbf{BNF-13}~: l'interface web doit être
    \textit{responsive} et compatible avec les navigateurs
    modernes. Elle est construite sur Bootstrap~5, dont le système
    de grille assure l'adaptation automatique de la mise en page
    aux différentes résolutions d'écran.

    \item \textbf{BNF-14}~: l'application est configurée
    principalement pour utiliser PostgreSQL
    (\texttt{django.db.backends.postgresql}) dans les environnements
    de déploiement. Une base SQLite peut également être utilisée en
    environnement local ou de développement afin de simplifier les
    phases de test et de prototypage.

    \item \textbf{BNF-15}~: le démarrage de l'ensemble de la pile
    applicative (service FastAPI sur le port~8000 et serveur Django
    sur le port~8080) est automatisé via le script PowerShell
    \texttt{start\_all.ps1}, permettant une mise en route
    reproductible en une seule commande.

\end{itemize}

\section{Modélisation des cas d'utilisation}

\subsection{Diagramme global des cas d'utilisation}

\begin{figure}[htbp]
    \centering
    \includegraphics[width=0.9\linewidth,height=9cm]{images/usercase.drawio.png}
    \caption{Diagramme global des cas d'utilisation}
    \label{fig:usercase_diag_global}
\end{figure}

Le diagramme global ci-dessus présente les interactions entre les
différents acteurs et les principaux cas d'utilisation du système.
Les acteurs de rang inférieur héritent des cas d'utilisation des
acteurs de rang supérieur selon la hiérarchie fonctionnelle définie
précédemment.

\subsection{Description détaillée des cas d'utilisation principaux}

\subsubsection{UC-01~— S'authentifier}

\begin{table}[htbp]
    \centering
    \caption{Description du cas d'utilisation UC-01~: S'authentifier}
    \label{tab:uc01}
    \begin{tabular}{|p{3.5cm}|p{10cm}|}
        \hline
        \textbf{Identifiant} & UC-01 \\
        \hline
        \textbf{Nom} & S'authentifier \\
        \hline
        \textbf{Acteurs} & Tous les utilisateurs \\
        \hline
        \textbf{Préconditions} & Le compte de l'utilisateur est
        actif (\texttt{status = active}). \\
        \hline
        \textbf{Scénario nominal} &
        \begin{enumerate}
            \item L'utilisateur accède à
            \texttt{/accounts/login/} et saisit son nom
            d'utilisateur et son mot de passe.
            \item Le système vérifie les credentials en base.
            \item En cas de succès, un code OTP à six chiffres
            est généré avec une expiration à dix minutes et
            transmis à l'adresse email de l'utilisateur via
            SendGrid.
            \item L'utilisateur est redirigé vers
            \texttt{/accounts/verify-otp/} et saisit le code
            reçu.
            \item Le système valide le code, marque l'OTP comme
            utilisé, ouvre la session et journalise la connexion
            (IP, \textit{user-agent}, horodatage).
            \item L'utilisateur est redirigé vers le tableau de
            bord principal (ou vers la page d'accueil restreinte
            s'il est un agent).
        \end{enumerate} \\
        \hline
        \textbf{Scénarios alternatifs} &
        \textbf{A1}~— Identifiants incorrects~: le système
        affiche un message d'erreur générique sans préciser
        lequel des deux champs est erroné. \\
        & \textbf{A2}~— Compte en attente de validation~: l'accès
        est refusé avec un message informant l'utilisateur que
        son compte n'a pas encore été activé. \\
        & \textbf{A3}~— Code OTP expiré ou invalide~: un message
        d'erreur est affiché et l'utilisateur peut en demander
        un nouveau. \\
        \hline
        \textbf{Postconditions} & Une session Django est ouverte.
        Un enregistrement \texttt{LoginActivity} est créé. \\
        \hline
    \end{tabular}
\end{table}

\newpage

\subsubsection{UC-02~— Lancer une analyse}

\begin{table}[htbp]
    \centering
    \caption{Description du cas d'utilisation UC-02~: Lancer une analyse}
    \label{tab:uc02}
    \begin{tabular}{|p{3.5cm}|p{10cm}|}
        \hline
        \textbf{Identifiant} & UC-02 \\
        \hline
        \textbf{Nom} & Lancer une analyse de portefeuille \\
        \hline
        \textbf{Acteurs} & Agent Marketing, Super Admin \\
        \hline
        \textbf{Préconditions} & L'utilisateur est authentifié et
        dispose du rôle \textit{agent\_marketing} ou
        \textit{super\_admin}. \\
        \hline
        \textbf{Scénario nominal} &
        \begin{enumerate}
            \item L'utilisateur accède à
            \texttt{/lancer-analyse/} et choisit une méthode
            d'alimentation~: import d'un fichier CSV ou
            génération de données mock.
            \item Le système crée un enregistrement
            \texttt{Dataset} traçant la source et charge les
            clients dans la table \texttt{ClientChurn}.
            \item Django interroge le service FastAPI via
            \texttt{GET /health}. Si le service est disponible,
            les données sont transmises à
            \texttt{POST /api/predict/batch}.
            \item FastAPI retourne les prédictions de churn et
            les scores associés pour l'ensemble du portefeuille.
            \item Les champs \texttt{score\_churn} et
            \texttt{churn\_predit} de chaque \texttt{ClientChurn}
            sont mis à jour.
            \item Une \texttt{AnalyseSession} est créée avec les
            métriques globales.
            \item Des notifications d'alerte in-app sont générées pour
            les clients identifiés comme étant à risque élevé.
        \end{enumerate} \\
        \hline
        \textbf{Postconditions} & Les scores des clients sont mis
        à jour. Une \texttt{AnalyseSession} est enregistrée. Des
        notifications in-app sont créées. \\
        \hline
    \end{tabular}
\end{table}

\subsubsection{UC-03~— Consulter la fiche d'un client}

\begin{table}[htbp]
    \centering
    \caption{Description du cas d'utilisation UC-03~: Consulter la fiche client}
    \label{tab:uc03}
    \begin{tabular}{|p{3.5cm}|p{10cm}|}
        \hline
        \textbf{Identifiant} & UC-03 \\
        \hline
        \textbf{Nom} & Consulter la fiche d'un client \\
        \hline
        \textbf{Acteurs} & Agent Commercial, Chef d'Agence, Agent
        Marketing \\
        \hline
        \textbf{Préconditions} & L'utilisateur est authentifié.
        Le client appartient à l'agence de l'utilisateur. \\
        \hline
        \textbf{Scénario nominal} &
        \begin{enumerate}
            \item L'utilisateur accède à la liste des clients
            (\texttt{/clients/}) et applique les filtres de
            recherche disponibles si nécessaire.
            \item Il sélectionne un client. Le système charge
            ses données depuis la base PostgreSQL.
            \item La fiche affiche le profil complet du client,
            son score de churn, sa prédiction binaire ainsi que
            les données d'interprétation SHAP disponibles,
            illustrant les variables les plus contributives à sa
            prédiction.
            \item Les recommandations existantes et leur statut
            sont affichés en bas de la fiche.
            \item L'utilisateur peut déclencher un export PDF de
            la fiche (\texttt{/clients/<id>/pdf/}).
        \end{enumerate} \\
        \hline
        \textbf{Postconditions} & La fiche est consultée. Un
        export PDF peut être téléchargé. \\
        \hline
    \end{tabular}
\end{table}

\subsubsection{UC-04~— Gérer une recommandation}

\begin{table}[htbp]
    \centering
    \caption{Description du cas d'utilisation UC-04~: Gérer une recommandation}
    \label{tab:uc04}
    \begin{tabular}{|p{3.5cm}|p{10cm}|}
        \hline
        \textbf{Identifiant} & UC-04 \\
        \hline
        \textbf{Nom} & Gérer une recommandation \\
        \hline
        \textbf{Acteurs} & Agent Commercial (création), Chef
        d'Agence (validation) \\
        \hline
        \textbf{Préconditions} & L'agent est authentifié avec le
        rôle \textit{agent\_commercial}. Le client cible est
        identifié dans son agence. \\
        \hline
        \textbf{Scénario nominal} &
        \begin{enumerate}
            \item L'agent rédige une recommandation depuis la
            fiche client et la soumet. Le statut est défini à
            \textit{pending}.
            \item Le système crée un enregistrement
            \texttt{Recommandation} et génère une notification
            in-app à destination du chef d'agence.
            \item Le chef d'agence consulte la recommandation et
            choisit de la valider ou de la rejeter.
            \item En cas de rejet, un motif est saisi et
            enregistré dans \texttt{RejetRecommandation}.
            \item Dans les deux cas, l'agent est notifié de
            l'issue par notification in-app uniquement.
        \end{enumerate} \\
        \hline
        \textbf{Postconditions} & La recommandation est en statut
        \textit{validated} ou \textit{rejected}. Les deux acteurs
        ont reçu leurs notifications in-app. \\
        \hline
    \end{tabular}
\end{table}

\section{Matrice de couverture fonctionnelle}

Le tableau suivant synthétise les fonctionnalités accessibles à
chaque profil d'utilisateur, reflétant ainsi la politique de
contrôle d'accès basée sur les rôles appliquée dans l'ensemble de
la plateforme.

\vspace{2cm}

\begin{table}[htbp]
    \centering
    \caption{Matrice de couverture fonctionnelle par rôle}
    \label{tab:matrice_roles}
    \begin{tabular}{|p{4.5cm}|c|c|c|c|c|}
        \hline
        \textbf{Fonctionnalité}
        & \rotatebox{70}{\textbf{Super Admin}}
        & \rotatebox{70}{\textbf{Admin Ville}}
        & \rotatebox{70}{\textbf{Chef Agence}}
        & \rotatebox{70}{\textbf{Ag. Commercial}}
        & \rotatebox{70}{\textbf{Ag. Marketing}} \\
        \hline
        Gestion des comptes       & \checkmark & \checkmark &  \checkmark          &            &            \\
        \hline
        Authentification 2FA      & \checkmark & \checkmark & \checkmark & \checkmark & \checkmark \\
        \hline
        Tableau de bord KPIs      & \checkmark & \checkmark & \checkmark &            &            \\
        \hline
        Liste clients             & \checkmark & \checkmark & \checkmark & \checkmark & \checkmark \\
        \hline
        Fiche client + SHAP       & \checkmark & \checkmark & \checkmark & \checkmark & \checkmark \\
        \hline
        Export PDF fiche client   & \checkmark & \checkmark & \checkmark & \checkmark & \checkmark \\
        \hline
        Créer recommandation      &            &            &            & \checkmark &    \checkmark        \\
        \hline
        Valider recommandation    & \checkmark &            & \checkmark &            &            \\
        \hline
        Importer dataset          & \checkmark &            & \checkmark &            &            \\
        \hline
        Lancer analyse            & \checkmark &            & \checkmark &            &            \\
        \hline
        Générer données mock      & \checkmark &            & \checkmark &            &            \\
        \hline
        Réinitialiser dashboard   & \checkmark &            &            &            &            \\
        \hline
        Administration système    & \checkmark &            &            &            &            \\
        \hline
    \end{tabular}
\end{table}

\section{Backlog produit et planification des sprints}
\label{sec:backlog}

Conformément à la méthodologie Scrum retenue pour le volet
applicatif et au processus CRISP-DM appliqué au volet data science,
l'ensemble des fonctionnalités du projet a été formalisé sous forme
de \textit{user stories} regroupées dans un \textbf{backlog produit}
unique. Ce backlog a ensuite été découpé en \textbf{six sprints}
successifs, dont les trois premiers correspondent aux phases data
science du projet et les trois derniers à la conception et au
développement de la plateforme web.

\subsection{Conventions de notation}

Chaque \textit{user story} est désignée par un identifiant unique
de la forme \texttt{US-XX} et caractérisée par~:

\begin{itemize}
    \item un \textbf{intitulé} formulé selon le canevas standard~:
    \og en tant que \textit{rôle}, je veux \textit{action}, afin
    de \textit{bénéfice} \fg{}~;
    \item une \textbf{priorité} sur trois niveaux~: \textbf{Haute},
    \textbf{M}oyenne, \textbf{B}asse, reflétant la criticité de la
    fonctionnalité pour la valeur métier du produit~;
    \item une \textbf{estimation d'effort} en \textit{Story Points}
    selon la suite de Fibonacci (1, 2, 3, 5, 8, 13), représentant
    la complexité relative et non un nombre d'heures absolu.
\end{itemize}

\subsection{Backlog produit consolidé}

Le tableau~\ref{tab:backlog_produit} présente l'intégralité du
backlog produit, organisé par épopée fonctionnelle
(\textit{Epic}). Cinquante \textit{user stories} composent ce
backlog, couvrant à la fois le pipeline data science et la
plateforme applicative.

\begin{footnotesize}
\begin{longtable}{|p{1.2cm}|p{8.5cm}|c|c|c|}
\caption{Backlog produit consolidé}
\label{tab:backlog_produit} \\
\hline
\textbf{ID} & \textbf{User Story} &
\textbf{Prio.} & \textbf{SP} & \textbf{Sprint} \\
\hline
\endfirsthead
\multicolumn{5}{c}{\textit{Suite du
tableau~\ref{tab:backlog_produit}}} \\
\hline
\textbf{ID} & \textbf{User Story} &
\textbf{Prio.} & \textbf{SP} & \textbf{Sprint} \\
\hline
\endhead
\hline
\endfoot

\multicolumn{5}{|l|}{\cellcolor{gray!15}\textbf{Epic 1~:
Construction du jeu de données synthétique}} \\
\hline
US-01 & En tant que data scientist, je veux concevoir un générateur
de profils clients tunisiens, afin de disposer d'un dataset
conforme au contexte de Tunisie Telecom.
& H & 8 & S1 \\
\hline
US-02 & En tant que data scientist, je veux générer les événements
CDR (appels, SMS, data), afin de modéliser l'usage télécom de
chaque client.
& H & 5 & S1 \\
\hline
US-03 & En tant que data scientist, je veux générer les sessions
digitales (empreinte numérique MyTT), afin de modéliser
l'engagement digital.
& H & 5 & S1 \\
\hline
US-04 & En tant que data scientist, je veux générer les points de
géolocalisation et la qualité de signal, afin de capter la
dimension territoriale du churn.
& M & 3 & S1 \\
\hline
US-05 & En tant que data scientist, je veux simuler des valeurs
manquantes selon les mécanismes MCAR et MAR, afin de rapprocher le
dataset des conditions réelles.
& H & 3 & S1 \\
\hline
US-06 & En tant que data scientist, je veux générer l'historique
des réclamations avec gradient de satisfaction, afin d'introduire
un signal d'insatisfaction.
& M & 3 & S1 \\
\hline
US-07 & En tant que data scientist, je veux définir et calculer la
variable cible churn selon la règle RGS-90, afin de disposer d'une
variable cible réglementairement ancrée.
& H & 5 & S1 \\
\hline
US-08 & En tant que data scientist, je veux conduire un audit
qualité du dataset (anomalies A, B, C, D), afin de garantir la
cohérence métier des données.
& H & 5 & S1 \\
\hline
US-09 & En tant que data scientist, je veux encapsuler toute la
logique dans un module \texttt{generator.py} paramétrable, afin de
pouvoir régénérer le dataset à la demande.
& M & 3 & S1 \\
\hline
US-10 & En tant que data scientist, je veux valider statistiquement
le dataset synthétique via SDV (KS, $\chi^2$), afin de m'assurer de
sa cohérence distributionnelle.
& M & 5 & S1 \\
\hline

\multicolumn{5}{|l|}{\cellcolor{gray!15}\textbf{Epic 2~: Analyse
exploratoire et ingénierie des features}} \\
\hline
US-11 & En tant que data scientist, je veux détecter la
multicolinéarité par le VIF, afin d'éliminer les variables
redondantes avant la modélisation.
& H & 3 & S2 \\
\hline
US-12 & En tant que data scientist, je veux calculer les
statistiques univariées conditionnelles, afin de comprendre la
structure du dataset par segment.
& M & 3 & S2 \\
\hline
US-13 & En tant que data scientist, je veux conduire l'analyse
bivariée par statut de churn, plan tarifaire et type d'abonnement,
afin d'identifier les segments à risque.
& H & 5 & S2 \\
\hline
US-14 & En tant que data scientist, je veux calculer la matrice de
corrélation de Spearman, afin d'identifier les relations monotones
entre features.
& H & 3 & S2 \\
\hline
US-15 & En tant que data scientist, je veux construire quatre
features composites (tendance data, ratios, score qualité zone,
score frustration), afin d'enrichir la capacité prédictive du
modèle.
& H & 5 & S2 \\
\hline
US-16 & En tant que data scientist, je veux produire les
visualisations exploratoires (countplots, histogrammes, boxplots,
scatter), afin de documenter visuellement le phénomène du churn.
& M & 5 & S2 \\
\hline
US-17 & En tant que data scientist, je veux arrêter définitivement
la matrice de 28 features, afin de figer l'entrée du pipeline de
modélisation.
& H & 2 & S2 \\
\hline

\multicolumn{5}{|l|}{\cellcolor{gray!15}\textbf{Epic 3~:
Modélisation prédictive et explicabilité}} \\
\hline
US-18 & En tant que data scientist, je veux préparer finalement les
données (encodage, split stratifié, imputation différenciée), afin
d'éviter toute fuite de données.
& H & 5 & S3 \\
\hline
US-19 & En tant que data scientist, je veux entraîner trois modèles
baseline (LR, RF, XGBoost) avec gestion du déséquilibre, afin de
disposer d'un point de comparaison.
& H & 5 & S3 \\
\hline
US-20 & En tant que data scientist, je veux optimiser les
hyperparamètres par GridSearchCV avec scoring $F_{\beta=2}$, afin
d'améliorer le rappel des modèles.
& H & 8 & S3 \\
\hline
US-21 & En tant que data scientist, je veux conduire une
optimisation bayésienne par Optuna, afin d'explorer plus finement
l'espace des hyperparamètres.
& H & 8 & S3 \\
\hline
US-22 & En tant que data scientist, je veux conduire une
analyse économique des modèles candidats, afin de retenir celui qui
minimise le coût métier total.
& H & 5 & S3 \\
\hline
US-23 & En tant que data scientist, je veux implémenter
l'explicabilité TreeSHAP (globale et individuelle), afin de
produire des prédictions interprétables par les agents.
& H & 5 & S3 \\
\hline
US-24 & En tant que data scientist, je veux sérialiser les
artefacts ML (modèle, seuil, features, explainer), afin de
permettre leur consommation par l'API.
& H & 3 & S3 \\
\hline

\multicolumn{5}{|l|}{\cellcolor{gray!15}\textbf{Epic 4~: API
d'inférence et sécurité applicative}} \\
\hline
US-25 & En tant qu'ingénieur backend, je veux concevoir l'API
FastAPI avec pattern Lifespan, afin de charger les artefacts une
seule fois au démarrage.
& H & 5 & S4 \\
\hline
US-26 & En tant qu'ingénieur backend, je veux définir les schémas
Pydantic d'entrée et de sortie, afin de garantir l'intégrité des
requêtes.
& H & 3 & S4 \\
\hline
US-27 & En tant qu'ingénieur backend, je veux implémenter les
endpoints \texttt{/health} et \texttt{/api/model/info}, afin de
superviser le service en production.
& M & 2 & S4 \\
\hline
US-28 & En tant qu'agent, je veux pouvoir obtenir une prédiction
individuelle expliquée pour un client, afin de comprendre le risque
de churn.
& H & 5 & S4 \\
\hline
US-29 & En tant qu'agent marketing, je veux pouvoir lancer une
prédiction par lot triée par risque, afin de prioriser mes actions
de rétention.
& H & 5 & S4 \\
\hline
US-30 & En tant qu'utilisateur, je veux pouvoir m'inscrire et faire
valider mon compte par mon supérieur hiérarchique, afin d'accéder à
la plateforme.
& H & 5 & S4 \\
\hline
US-31 & En tant qu'utilisateur, je veux m'authentifier en double
facteur via OTP par email, afin de sécuriser l'accès à mon compte.
& H & 8 & S4 \\
\hline
US-32 & En tant qu'utilisateur, je veux pouvoir réinitialiser mon
mot de passe oublié, afin de récupérer l'accès à mon compte.
& M & 3 & S4 \\
\hline
US-33 & En tant qu'administrateur, je veux que les endpoints
sensibles soient protégés contre les attaques par force brute, afin
de sécuriser la plateforme.
& H & 3 & S4 \\
\hline
US-34 & En tant qu'administrateur, je veux que toutes les
connexions soient journalisées (IP, user-agent, horodatage), afin
de disposer d'un journal d'audit.
& M & 2 & S4 \\
\hline

\multicolumn{5}{|l|}{\cellcolor{gray!15}\textbf{Epic 5~: Données
clients et intégration ML}} \\
\hline
US-35 & En tant que chef d'agence, je veux pouvoir importer un
dataset client au format CSV, afin d'alimenter la plateforme.
& H & 5 & S5 \\
\hline
US-36 & En tant que chef d'agence, je veux que la structure du
fichier soit validée automatiquement, afin d'être alerté
immédiatement en cas d'anomalie.
& H & 3 & S5 \\
\hline
US-37 & En tant que chef d'agence, je veux pouvoir générer un
dataset de clients simulés, afin de tester la plateforme sans
donnée réelle.
& M & 3 & S5 \\
\hline
US-38 & En tant que chef d'agence, je veux pouvoir lancer une
analyse de mon portefeuille, afin de produire les scores de churn
pour tous mes clients.
& H & 5 & S5 \\
\hline
US-39 & En tant qu'utilisateur, je veux que mon analyse bascule
automatiquement sur le pipeline local si FastAPI est indisponible,
afin d'assurer la continuité de service.
& H & 5 & S5 \\
\hline
US-40 & En tant qu'utilisateur, je veux que les données de mon
agence soient strictement isolées des autres, afin de garantir la
confidentialité.
& H & 5 & S5 \\
\hline
US-41 & En tant qu'utilisateur, je veux que chaque analyse soit
historisée avec ses métriques globales, afin de suivre l'évolution
dans le temps.
& M & 3 & S5 \\
\hline

\multicolumn{5}{|l|}{\cellcolor{gray!15}\textbf{Epic 6~: Interface
utilisateur et fonctionnalités métier}} \\
\hline
US-42 & En tant qu'utilisateur, je veux disposer d'un tableau de
bord avec les KPIs clés du portefeuille, afin d'avoir une vision
synthétique en un coup d'\oe il. \textit{(Réservé aux chefs d'agence et administrateurs)}
& H & 5 & S6 \\
\hline
US-43 & En tant qu'utilisateur, je veux pouvoir consulter la liste
paginée des clients avec filtres par risque, afin de naviguer
efficacement dans mon portefeuille.
& H & 3 & S6 \\
\hline
US-44 & En tant qu'utilisateur, je veux consulter la fiche
détaillée d'un client avec son explication SHAP, afin de comprendre
les facteurs de son risque.
& H & 5 & S6 \\
\hline
US-45 & En tant qu'agent commercial, je veux pouvoir créer une
recommandation pour un client à risque, afin de proposer une action
de rétention.
& H & 3 & S6 \\
\hline
US-46 & En tant que chef d'agence, je veux pouvoir valider ou
rejeter les recommandations soumises par mes agents, afin de garder
le contrôle sur les actions menées.
& H & 5 & S6 \\
\hline
US-47 & En tant qu'utilisateur, je veux recevoir des notifications
in-app pour les événements métier, afin d'être alerté en temps réel.
\textit{(Les notifications par email sont réservées aux événements liés au compte et à l'authentification)}
& H & 5 & S6 \\
\hline
US-48 & En tant qu'utilisateur, je veux pouvoir gérer mes
notifications (lire, archiver, supprimer), afin de maintenir un
centre de notifications propre.
& M & 2 & S6 \\
\hline
US-49 & En tant qu'utilisateur, je veux pouvoir exporter la fiche
d'un client en PDF, afin de la partager ou de l'archiver hors
plateforme.
& M & 3 & S6 \\
\hline
US-50 & En tant qu'administrateur, je veux que l'interface soit
responsive et adaptée à tout type d'écran, afin que la plateforme
soit utilisable depuis n'importe quel terminal.
& M & 3 & S6 \\
\hline
\end{longtable}
\end{footnotesize}

\subsection{Récapitulatif des sprints}

Les cinquante \textit{user stories} du backlog ont été réparties
sur six sprints successifs d'une durée moyenne de deux à trois
semaines. Le tableau~\ref{tab:recap_sprints} en présente la
synthèse.

\begin{table}[htbp]
\centering
\caption{Synthèse des six sprints du projet}
\label{tab:recap_sprints}
\small
\begin{tabular}{|c|p{5.5cm}|c|c|p{3.5cm}|}
\hline
\textbf{Sprint} & \textbf{Objectif principal}
& \textbf{Stories} & \textbf{Total SP}
& \textbf{Livrable clé} \\
\hline
S1 & Construction du jeu de données synthétique
& 10 & 45 & Dataset 300 clients validé \\
\hline
S2 & Analyse exploratoire et sélection des features
& 7 & 26 & Matrice de 28 features \\
\hline
S3 & Modélisation prédictive et explicabilité
& 7 & 39 & Modèle RF Optuna sérialisé \\
\hline
S4 & API FastAPI et sécurité applicative
& 10 & 41 & Service FastAPI et auth OTP \\
\hline
S5 & Données clients et intégration ML
& 7 & 29 & Pipeline d'import et d'analyse \\
\hline
S6 & Interface utilisateur et fonctionnalités métier
& 9 & 34 & Plateforme web complète \\
\hline
\multicolumn{2}{|r|}{\textbf{Total}}
& \textbf{50} & \textbf{214} & \\
\hline
\end{tabular}
\end{table}

Cette planification reflète la nature hybride du projet~: les
trois premiers sprints couvrent intégralement le cycle CRISP-DM
(compréhension, préparation, modélisation, évaluation), tandis que
les trois suivants déploient les résultats au sein d'une plateforme
applicative complète. Les détails de chaque sprint (objectif,
\textit{Definition of Done} et stories retenues) sont présentés au
fil des chapitres correspondants.

\section{Conclusion}

Ce chapitre a permis de formaliser avec précision l'ensemble des
exigences auxquelles la plateforme devra répondre. L'identification
des acteurs et de leurs périmètres respectifs a mis en évidence la
complexité de la politique de contrôle d'accès à mettre en œuvre,
tandis que la formalisation des besoins fonctionnels a révélé
l'étendue des interactions entre les différents modules du système.

La distinction entre les besoins fonctionnels et non fonctionnels
s'avère particulièrement structurante~: les contraintes de
sécurité (authentification double facteur, isolation par agence,
limitation du nombre de requêtes) et les exigences de robustesse
(mécanisme de \textit{fallback} ML, continuité de service)
constituent des exigences de premier ordre qui influenceront
directement les choix architecturaux et d'implémentation.

L'ensemble de ces besoins a été traduit en un backlog produit de
cinquante \textit{user stories} réparties sur six sprints, qui
guideront l'organisation du travail tout au long du projet. Le
chapitre suivant amorce la première phase opérationnelle~: la
construction du jeu de données synthétique, qui constitue le
socle de tout le pipeline prédictif.
